% !TeX root = ../数学分析教材（伍胜健）笔记.tex
% 本文件由主文档通过 \input 载入；请编译主文档。
\setcounter{chapter}{16} % 保留原书章号。
\chapter{含参变量积分}
\label{chap:parameter-dependent-integrals}

% 本章依据“数学分析/textbooks/数学分析（分章原文）/数学分析-17.pdf”第 1--28 页习题十七之前的正文整理。
积分中除积分变量外还可以含有参数。例如，随时间变化的向量场通过一个固定曲面的瞬时通量，是对空间变量积分后得到的时间函数；求一段时间内的总通量，还需再对时间积分。本章研究的核心是：积分得到的参数函数是否连续、可导，以及极限、微分、积分能否与原积分交换。有限区间上的联合连续性提供一致控制；积分区间无穷或被积函数有瑕点时，还必须控制截断误差。本章以黎曼积分及相应广义积分为基础，先讨论一般理论，再用它研究 $\Gamma$ 函数与 $B$ 函数。

\section{含参变量定积分}
\label{sec:parameter-dependent-definite-integrals}

设 $a<b$、$c<d$，函数 $f(x,y)$ 定义在矩形 $D=[a,b]\times[c,d]$ 上。若对每个固定的 $x\in[a,b]$，$f(x,\cdot)$ 在 $[c,d]$ 上可积，则
\[
    I(x)=\int_c^d f(x,y)\,\mathrm{d}y
\]
定义了一个以 $x$ 为自变量的函数，称为\emph{含参变量定积分}（parameter-dependent definite integral）。这里 $y$ 是积分变量，$x$ 是参变量；交换二者的角色也得到同类问题。以下端点处的连续性与导数按相应单侧理解。

\begin{theorem}[连续性与极限交换]
\label{thm:parameter-finite-integral-continuity}
若 $f\in C(D)$，则 $I\in C([a,b])$。特别地，对任意 $x_0\in[a,b]$，
\begin{equation}
    \lim_{x\to x_0}\int_c^d f(x,y)\,\mathrm{d}y
    =\int_c^d f(x_0,y)\,\mathrm{d}y
    =\int_c^d\lim_{x\to x_0}f(x,y)\,\mathrm{d}y.
    \label{eq:parameter-finite-limit-exchange}
\end{equation}
\end{theorem}

\begin{proof}
$D$ 紧且 $f$ 连续，故 $f$ 在 $D$ 上一致连续。当 $x\to x_0$ 时，$\sup_{y\in[c,d]}|f(x,y)-f(x_0,y)|\to0$，于是
\[
    |I(x)-I(x_0)|\leq(d-c)
    \sup_{y\in[c,d]}|f(x,y)-f(x_0,y)|\longrightarrow0.
\]
这里将逐点连续性加强为对所有积分变量统一有效的估计，正是能够交换极限与积分的原因。
\end{proof}

同样，$F(x,u)=\int_c^u f(x,y)\,\mathrm{d}y$ 在 $D$ 上联合连续。事实上，若 $M=\max_D|f|$，则
\[
    |F(x,u)-F(x_0,u_0)|
    \leq(d-c)\sup_{y\in[c,d]}|f(x,y)-f(x_0,y)|+M|u-u_0|.
\]
联合连续只是积分函数连续的充分条件。例如 $f(x,y)=h(x)g(y)$，若 $g$ 可积且 $\int_c^d g=0$，则无论 $h$ 是否连续，都有 $I(x)=0$。

\begin{theorem}[积分次序交换]
\label{thm:parameter-finite-integration-exchange}
若 $f\in C(D)$，则两个含参变量定积分分别是相应闭区间上的连续函数，且
\begin{equation}
    \int_a^b\left(\int_c^d f(x,y)\,\mathrm{d}y\right)\mathrm{d}x
    =\int_c^d\left(\int_a^b f(x,y)\,\mathrm{d}x\right)\mathrm{d}y.
    \label{eq:parameter-finite-integration-exchange}
\end{equation}
\end{theorem}

\begin{proof}
连续性由定理~\ref{thm:parameter-finite-integral-continuity} 得到。连续函数在矩形上重可积，两个累次积分都等于 $\iint_D f(x,y)\,\mathrm{d}x\,\mathrm{d}y$，故相等。
\end{proof}

\begin{theorem}[积分号下求导]
\label{thm:parameter-finite-differentiation}
若 $f$ 与偏导数 $f_x=\partial f/\partial x$ 在 $D$ 上连续，则 $I\in C^1([a,b])$，且
\begin{equation}
    I'(x)=\frac{\mathrm{d}}{\mathrm{d}x}\int_c^d f(x,y)\,\mathrm{d}y
    =\int_c^d f_x(x,y)\,\mathrm{d}y.
    \label{eq:parameter-finite-differentiation}
\end{equation}
\end{theorem}

\begin{proof}
令 $g(x)=\int_c^d f_x(x,y)\,\mathrm{d}y$，由连续性定理知 $g$ 连续。对固定的 $y$ 使用微积分基本定理，再交换有限区间上的积分，得
\[
    I(x)-I(a)=\int_c^d\int_a^x f_x(t,y)\,\mathrm{d}t\,\mathrm{d}y
    =\int_a^x g(t)\,\mathrm{d}t.
\]
两端对 $x$ 求导，便得 $I'=g$。
\end{proof}

\begin{theorem}[变限积分的莱布尼茨公式]
\label{thm:parameter-leibniz-variable-limits}
设 $f,f_x\in C(D)$，$\alpha,\beta:[a,b]\to[c,d]$ 可导，则
\begin{equation}
    \frac{\mathrm{d}}{\mathrm{d}x}
      \int_{\alpha(x)}^{\beta(x)}f(x,y)\,\mathrm{d}y
    =\int_{\alpha(x)}^{\beta(x)}f_x(x,y)\,\mathrm{d}y
      +f(x,\beta(x))\beta'(x)-f(x,\alpha(x))\alpha'(x).
    \label{eq:parameter-leibniz-variable-limits}
\end{equation}
积分按定向区间理解，不要求始终有 $\alpha(x)\leq\beta(x)$。
\end{theorem}

\begin{proof}
对 $F(x,u)=\int_c^u f(x,y)\,\mathrm{d}y$，前一定理与微积分基本定理分别给出
\[
    F_x(x,u)=\int_c^u f_x(x,y)\,\mathrm{d}y,
    \qquad F_u(x,u)=f(x,u).
\]
这两个偏导数联合连续，故 $F$ 可微。把原积分写成 $F(x,\beta(x))-F(x,\alpha(x))$，用复合函数求导法则即得结论。若积分限在端点取到 $c$ 或 $d$，相应增量沿矩形内取值，仍有同样的单侧微分估计。
\end{proof}

\emph[tip]{变限求导包含两类贡献：被积函数随参数的变化，以及上下限移动所增减的积分。}上下限固定时只保留积分项；被积函数不显含 $x$ 时只保留边界项。

\begin{example}
\label{ex:parameter-finite-limit-and-logarithm}
求 $\displaystyle\lim_{x\to0+}\int_0^1(1+y)^x e^{xy}\sin(x+y)\,\mathrm{d}y$，并计算
\[
    \int_0^1\frac{x^{2e-1}-x^{e-1}}{\ln x}\,\mathrm{d}x.
\]
\end{example}

\begin{solve}
第一个被积函数在 $[0,1]\times[0,1]$ 上连续，故极限等于
\[
    \int_0^1\sin y\,\mathrm{d}y=1-\cos1.
\]
第二个积分利用参数积分表示
\[
    \frac{x^{2e-1}-x^{e-1}}{\ln x}
    =\int_{e-1}^{2e-1}x^y\,\mathrm{d}y\qquad(0<x<1).
\]
因为 $y\geq e-1>0$，把 $x^y$ 在 $x=0$ 处补为零后，它在 $[0,1]\times[e-1,2e-1]$ 上联合连续；在 $x=1$ 处原商的连续延拓值为 $e$。因此交换积分次序合法，并有
\[
    \int_{e-1}^{2e-1}\int_0^1 x^y\,\mathrm{d}x\,\mathrm{d}y
    =\int_{e-1}^{2e-1}\frac{\mathrm{d}y}{1+y}
    =\ln2.
\]
参数积分把原来不易直接求原函数的商，化成了幂函数的简单积分。
\end{solve}

\begin{example}
\label{ex:parameter-nested-variable-limits}
设 $f\in C(\mathbb R^2)$，令
\[
    F(x)=\int_0^{x^2}\left(\int_0^x f(s,t)\,\mathrm{d}t\right)\mathrm{d}s.
\]
求 $F'(x)$。
\end{example}

\begin{solve}
记 $g(s,x)=\int_0^x f(s,t)\,\mathrm{d}t$，则 $g$ 联合连续，且 $g_x(s,x)=f(s,x)$ 联合连续。对任意固定 $x$，在包含所有相关积分限的有限矩形上应用式~\eqref{eq:parameter-leibniz-variable-limits}，得
\[
    F'(x)=\int_0^{x^2}f(s,x)\,\mathrm{d}s
      +2x\int_0^x f(x^2,t)\,\mathrm{d}t.
\]
这里只有 $f$ 的连续性，不需要 $f$ 本身有偏导数；所需的参数导数来自内层变上限积分。
\end{solve}

\section{含参变量广义积分}
\label{sec:parameter-dependent-improper-integrals}

把有限积分区间改为无穷区间，或允许被积函数在端点无界后，积分还包含一个截断极限。因此，有限截断积分的连续性或可导性，只解决了截断之前的问题；要把性质传给广义积分，还需控制参数变化时的尾部误差。

\subsection{含参变量无穷积分}
\label{subsec:parameter-dependent-infinite-integrals}

以下设 $E\subseteq\mathbb R$ 非空，$f$ 定义在 $E\times[c,+\infty)$ 上，且对每个 $x\in E$，$f(x,\cdot)$ 在每个有限区间 $[c,A]$ 上黎曼可积。若
\[
    I(x)=\int_c^{+\infty}f(x,y)\,\mathrm{d}y
    =\lim_{A\to+\infty}F_A(x),\qquad
    F_A(x)=\int_c^A f(x,y)\,\mathrm{d}y
\]
对每个 $x\in E$ 均存在且有限，就得到一个\emph{含参变量无穷积分}（parameter-dependent improper integral over an infinite interval）。

\begin{definition}[一致收敛与绝对一致收敛]
\label{def:parameter-infinite-uniform-convergence}
若上述积分在 $E$ 上逐点收敛，且
\begin{equation}
    \forall\varepsilon>0\ \exists A_0>c\ \forall A>A_0\ \forall x\in E,
    \quad\left|\int_A^{+\infty}f(x,y)\,\mathrm{d}y\right|<\varepsilon,
    \label{eq:parameter-infinite-uniform-tail}
\end{equation}
则称它在 $E$ 上\emph{一致收敛}（uniform convergence）。若 $\int_c^{+\infty}|f(x,y)|\,\mathrm{d}y$ 在 $E$ 上一致收敛，则称原积分在 $E$ 上\emph{绝对一致收敛}（uniform absolute convergence）。
\end{definition}

逐点收敛允许截断阈值 $A_0$ 随 $x$ 改变，一致收敛要求同一个阈值适用于所有参数。绝对一致收敛蕴含一致收敛，也蕴含每个参数处的绝对收敛；每个参数处绝对收敛本身不能保证一致收敛。与函数序列不同，截断参数 $A$ 是连续变化的实数，不能未经论证就只检验整数截断。

\begin{theorem}[一致收敛的柯西准则]
\label{thm:parameter-infinite-uniform-cauchy}
含参变量无穷积分在 $E$ 上一致收敛，当且仅当
\begin{equation}
    \forall\varepsilon>0\ \exists A_0>c\ \forall B>A>A_0\ \forall x\in E,
    \quad\left|\int_A^B f(x,y)\,\mathrm{d}y\right|<\varepsilon.
    \label{eq:parameter-infinite-uniform-cauchy}
\end{equation}
\end{theorem}

\begin{proof}
必要性由 $\int_A^B f=\int_A^{+\infty}f-\int_B^{+\infty}f$ 及三角不等式得到。反之，固定 $x$ 后，式~\eqref{eq:parameter-infinite-uniform-cauchy} 是普通无穷积分的柯西条件，故 $I(x)$ 存在。在该式中先把误差取为 $\varepsilon/2$，再令 $B\to+\infty$，可得对所有 $x$ 统一成立的尾积分估计 $|\int_A^{+\infty}f|\leq\varepsilon/2<\varepsilon$。
\end{proof}

\begin{theorem}[魏尔斯特拉斯判别法]
\label{thm:parameter-infinite-weierstrass}
若存在与 $x$ 无关的非负局部可积函数 $g$，使
\[
    |f(x,y)|\leq g(y)\quad(x\in E,\ y\geq c),\qquad
    \int_c^{+\infty}g(y)\,\mathrm{d}y<+\infty,
\]
则 $\int_c^{+\infty}f(x,y)\,\mathrm{d}y$ 在 $E$ 上绝对一致收敛。
\end{theorem}

\begin{proof}
对任意 $B>A\geq c$ 与 $x\in E$，$\int_A^B|f(x,y)|\,\mathrm{d}y\leq\int_A^B g(y)\,\mathrm{d}y$，右端满足与 $x$ 无关的柯西估计，应用前一定理即可。
\end{proof}

\begin{theorem}[狄利克雷与阿贝尔判别法]
\label{thm:parameter-infinite-dirichlet-abel}
设 $f(x,\cdot),g(x,\cdot)$ 在每个有限截断区间上可积，并且对每个固定 $x\in E$，$g(x,y)$ 关于 $y$ 单调。以下各组条件均能保证 $\int_c^{+\infty}f(x,y)g(x,y)\,\mathrm{d}y$ 在 $E$ 上一致收敛。

狄利克雷条件（Dirichlet test）：存在 $M>0$，使 $|\int_c^A f(x,y)\,\mathrm{d}y|\leq M$ 对所有 $A>c$、$x\in E$ 成立，且 $g(x,y)$ 在 $y\to+\infty$ 时关于 $x\in E$ 一致趋于零。阿贝尔条件（Abel test）：$\int_c^{+\infty}f(x,y)\,\mathrm{d}y$ 在 $E$ 上一致收敛，且存在 $L>0$，使 $|g(x,y)|\leq L$ 对所有 $x\in E$、$y\geq c$ 成立。
\end{theorem}

\begin{proof}
由定积分第二中值定理，对任意 $B>A\geq c$，存在 $\xi\in[A,B]$，使
\[
    \int_A^B fg\,\mathrm{d}y
    =g(x,A)\int_A^\xi f\,\mathrm{d}y
      +g(x,B)\int_\xi^B f\,\mathrm{d}y.
\]
在狄利克雷条件下，右侧两个 $f$ 的积分绝对值均不超过 $2M$，故原积分绝对值不超过 $2M(|g(x,A)|+|g(x,B)|)$，并一致趋于零。在阿贝尔条件下，一致柯西准则保证右侧两个积分对足够大的 $A$ 均小于任意指定的 $\delta>0$，故乘积积分绝对值不超过 $2L\delta$。两种情形都满足式~\eqref{eq:parameter-infinite-uniform-cauchy}。
\end{proof}

\begin{example}
\label{ex:parameter-exponential-tail-domain}
讨论 $\displaystyle\int_1^{+\infty}y^x e^{-y}\,\mathrm{d}y$ 在 $[0,M]$（$0<M<+\infty$）及 $[0,+\infty)$ 上的一致收敛性。
\end{example}

\begin{solve}
在 $[0,M]$ 上，因 $y\geq1$，有 $0\leq y^x e^{-y}\leq y^M e^{-y}$。后者可积，因为充分大的 $y$ 满足 $y^M\leq e^{y/2}$，故由魏尔斯特拉斯判别法绝对一致收敛。在 $[0,+\infty)$ 上，任给 $A>1$，都有
\[
    \int_A^{A+1}y^x e^{-y}\,\mathrm{d}y
    \geq A^x e^{-(A+1)}.
\]
选取足够大的 $x$ 可使右端至少为 $1$，所以一致柯西条件失败。每个固定 $x$ 处的绝对收敛，并不能抵消参数无界增长造成的尾部失控。
\end{solve}

\begin{example}
\label{ex:parameter-sine-tail-uniformity}
证明 $\displaystyle\int_1^{+\infty}\frac{\sin(xy)}{y}\,\mathrm{d}y$ 在每个 $[a_0,+\infty)$（$a_0>0$）上一致收敛，但在 $(0,+\infty)$ 上不一致收敛。
\end{example}

\begin{solve}
当 $x\geq a_0$ 时，$|\int_1^A\sin(xy)\,\mathrm{d}y|\leq2/a_0$，而 $1/y$ 单调趋于零且与参数无关，故一致狄利克雷判别法适用。为证后一结论，取 $A_n=n$、$B_n=2n$、$x_n=1/n$，换元 $t=y/n$ 得
\[
    \int_{A_n}^{B_n}\frac{\sin(x_ny)}{y}\,\mathrm{d}y
    =\int_1^2\frac{\sin t}{t}\,\mathrm{d}t>0.
\]
右端是不随 $n$ 改变的正常数，故尾积分不能一致趋于零。参数靠近零时，振荡抵消所需的积分长度不断增大。
\end{solve}

\begin{example}
\label{ex:parameter-abel-damping-equivalence}
设 $h\in C([0,+\infty))$。证明 $\int_0^{+\infty}e^{-\alpha t}h(t)\,\mathrm{d}t$ 对 $\alpha\in(0,+\infty)$ 一致收敛，当且仅当 $\int_0^{+\infty}h(t)\,\mathrm{d}t$ 收敛。
\end{example}

\begin{solve}
若无阻尼积分收敛，把 $h(t)$ 看成与 $\alpha$ 无关的函数，则它的积分关于 $\alpha$ 一致收敛；$e^{-\alpha t}$ 关于 $t$ 单调且统一有界于 $1$，阿贝尔判别法给出充分性。反之，若带阻尼积分一致收敛，则对任意 $\varepsilon>0$，存在 $A_0$，使所有 $B>A>A_0$ 及 $\alpha>0$ 都满足
\[
    \left|\int_A^B e^{-\alpha t}h(t)\,\mathrm{d}t\right|<\varepsilon/2.
\]
先固定有限的 $A,B$，再由有限区间积分的连续性令 $\alpha\to0+$，得到 $|\int_A^B h(t)\,\mathrm{d}t|\leq\varepsilon/2<\varepsilon$。普通无穷积分的柯西准则给出必要性。
\end{solve}

\subsection{含参变量无穷积分的性质}
\label{subsec:properties-of-parameter-dependent-infinite-integrals}

\begin{theorem}[连续截断与序列截断的联系]
\label{thm:parameter-infinite-sequential-criterion}
在局部可积的基本假设下，$\int_c^{+\infty}f(x,y)\,\mathrm{d}y$ 在 $E$ 上一致收敛，当且仅当对于每个严格递增且趋于 $+\infty$ 的实数序列 $\{t_n\}$（$t_1>c$），函数序列
\[
    F_{t_n}(x)=\int_c^{t_n}f(x,y)\,\mathrm{d}y
\]
在 $E$ 上一致收敛。此时各序列的极限都是同一个积分函数 $I(x)$。
\end{theorem}

\begin{proof}
必要性由定义直接得到。若积分不一致收敛，则一致柯西条件失败，存在 $\varepsilon_0>0$，可以逐次选择
\[
    c<A_1<B_1<A_2<B_2<\cdots,\qquad A_n>n,
\]
以及 $x_n\in E$，使 $|\int_{A_n}^{B_n}f(x_n,y)\,\mathrm{d}y|\geq\varepsilon_0$。将所有 $A_n,B_n$ 依次排成 $\{t_n\}$，所得截断函数序列不满足一致柯西准则，与假设矛盾。
\end{proof}

\emph[caution]{这里要求每个趋于无穷的递增截断序列都适用。}例如与参数无关的 $f(y)=\sin(2\pi y)$ 在整数截断处的积分恒为零，但 $\int_0^A f(y)\,\mathrm{d}y=(1-\cos(2\pi A))/(2\pi)$ 在 $A\to+\infty$ 时无极限。单独检验一列截断值会漏掉相邻截断之间的变化。

\begin{theorem}[连续性与有限区间上的积分交换]
\label{thm:parameter-infinite-continuity-integration}
设 $E$ 为区间，$f\in C(E\times[c,+\infty))$。若 $I(x)=\int_c^{+\infty}f(x,y)\,\mathrm{d}y$ 在 $E$ 上一致收敛，则 $I$ 在 $E$ 上连续。若 $E=[a,b]$，还有
\begin{equation}
    \int_a^b\left(\int_c^{+\infty}f(x,y)\,\mathrm{d}y\right)\mathrm{d}x
    =\int_c^{+\infty}\left(\int_a^b f(x,y)\,\mathrm{d}x\right)\mathrm{d}y,
    \label{eq:parameter-infinite-integration-exchange}
\end{equation}
其中右侧无穷积分也收敛。
\end{theorem}

\begin{proof}
每个有限截断积分 $F_A(x)$ 在 $E$ 上连续：对任意参数点，只需在其有限闭邻域与 $[c,A]$ 的乘积上应用定积分连续性定理。由一致极限保持连续性，$I$ 连续。在 $E=[a,b]$ 时，
\[
    \left|\int_a^b I(x)\,\mathrm{d}x-\int_a^bF_A(x)\,\mathrm{d}x\right|
    \leq(b-a)\sup_{x\in[a,b]}|I(x)-F_A(x)|\longrightarrow0.
\]
对 $F_A$ 的积分使用有限矩形上的次序交换，再令实数 $A\to+\infty$，便得到右侧广义积分的存在性及公式。
\end{proof}

若只要求 $I$ 在开区间内部连续，只需对每个闭子区间证明一致收敛，即\emph{内闭一致收敛}（locally uniform convergence）。但式~\eqref{eq:parameter-infinite-integration-exchange} 在整个 $[a,b]$ 上的结论，不能仅凭 $(a,b)$ 内部的局部一致收敛推出。

\begin{theorem}[无穷积分的积分号下求导]
\label{thm:parameter-infinite-differentiation}
设 $E$ 为非退化区间，$f,f_x\in C(E\times[c,+\infty))$。假设某个 $x_0\in E$ 处的积分 $\int_c^{+\infty}f(x_0,y)\,\mathrm{d}y$ 收敛，并且 $\int_c^{+\infty}f_x(x,y)\,\mathrm{d}y$ 在 $E$ 的每个闭子区间上一致收敛，则原积分在 $E$ 上处处收敛且内闭一致收敛，积分函数 $I$ 连续可导，并有
\begin{equation}
    I'(x)=\int_c^{+\infty}f_x(x,y)\,\mathrm{d}y.
    \label{eq:parameter-infinite-differentiation}
\end{equation}
若 $E$ 本身有界，且导数积分在整个 $E$ 上一致收敛，则原积分在整个 $E$ 上也一致收敛。
\end{theorem}

\begin{proof}
取任意闭子区间 $K\subseteq E$，必要时把它扩大为仍在 $E$ 内且包含 $x_0$ 的有限闭区间 $K_0$。由有限截断积分的求导公式与微积分基本定理，对 $B>A>c$ 及 $x\in K_0$，有
\[
    \int_A^B f(x,y)\,\mathrm{d}y
    =\int_A^B f(x_0,y)\,\mathrm{d}y
      +\int_{x_0}^x\left(\int_A^B f_x(t,y)\,\mathrm{d}y\right)\mathrm{d}t.
\]
因此其绝对值不超过
\[
    \left|\int_A^B f(x_0,y)\,\mathrm{d}y\right|
    +|K_0|\sup_{t\in K_0}
       \left|\int_A^B f_x(t,y)\,\mathrm{d}y\right|,
\]
其中 $|K_0|$ 表示区间长度。两项均趋于零，故原积分在 $K_0$ 上一致收敛。再任取 $t_n\uparrow+\infty$，截断函数 $F_{t_n}$ 在 $x_0$ 处收敛，其导函数在 $K_0$ 上一致收敛。由定理~\ref{thm:differentiation-of-function-sequence}，其极限可导且导数等于导函数的极限，即式~\eqref{eq:parameter-infinite-differentiation}。导数的连续性由前一连续性定理得到。若 $E$ 有界且导数积分在 $E$ 上一致收敛，上述估计中可用 $E$ 的有限长度统一控制 $|x-x_0|$，从而原积分在整个 $E$ 上一致收敛。
\end{proof}

\begin{remark}
\label{rem:parameter-derivative-unbounded-domain}
求导公式是局部结论，不能在无界参数区间上无条件加强为原积分整体一致收敛。例如 $f(x,y)=xe^{-y}$，$x\geq0$、$y\geq0$，则 $\int_0^{+\infty}f_x(x,y)\,\mathrm{d}y=1$ 关于 $x$ 一致收敛，$x=0$ 处原积分也收敛；但原积分的尾部为 $xe^{-A}$，在 $[0,+\infty)$ 上没有统一趋于零。此时 $I(x)=x$ 仍然可导，求导公式完全成立。
\end{remark}

\begin{proposition}[非负积分的狄尼定理]
\label{prop:parameter-integral-dini}
设 $f\in C([a,b]\times[c,+\infty))$，且 $f\geq0$。若 $I(x)=\int_c^{+\infty}f(x,y)\,\mathrm{d}y$ 对每个 $x\in[a,b]$ 收敛并在 $[a,b]$ 上连续，则该积分在 $[a,b]$ 上一致收敛。$f\leq0$ 时结论相同。
\end{proposition}

\begin{proof}
对任意 $t_n\uparrow+\infty$，连续函数 $F_{t_n}$ 单调增加并逐点趋于连续函数 $I$。由函数序列的狄尼定理，即定理~\ref{thm:dini-uniform-convergence}，它一致收敛。也可固定一列 $\{t_n\}$，由 $A\geq t_n$ 时 $0\leq I(x)-F_A(x)\leq I(x)-F_{t_n}(x)$，直接得到所有实数截断的一致估计。非正情形对 $-f$ 应用同一结论。
\end{proof}

\begin{example}
\label{ex:parameter-dirichlet-integral-evaluation}
利用含参变量积分计算狄利克雷积分 $\displaystyle\int_0^{+\infty}\frac{\sin t}{t}\,\mathrm{d}t$。
\end{example}

\begin{solve}
先确认积分存在：$\sin t/t$ 在零点可连续延拓为 $1$，而无穷尾部由普通狄利克雷判别法收敛。引入阻尼参数
\[
    I(\alpha)=\int_0^{+\infty}e^{-\alpha t}\frac{\sin t}{t}\,\mathrm{d}t,
    \qquad\alpha\geq0.
\]
无阻尼积分收敛，且 $e^{-\alpha t}$ 对 $t$ 单调、关于 $\alpha\geq0$ 统一有界于 $1$，故一致阿贝尔判别法保证 $I$ 的积分在 $[0,+\infty)$ 上一致收敛。被积函数在 $t=0$ 处按连续延拓理解，因此 $I$ 在 $\alpha=0$ 处右连续。

对任意 $\delta>0$，当 $\alpha\geq\delta$ 时，参数偏导数为 $-e^{-\alpha t}\sin t$，绝对值不超过可积函数 $e^{-\delta t}$。故对所有 $\alpha>0$ 可在积分号下求导，得到
\[
    I'(\alpha)=-\int_0^{+\infty}e^{-\alpha t}\sin t\,\mathrm{d}t
    =-\frac1{1+\alpha^2}.
\]
最后一步可用原函数 $-e^{-\alpha t}(\alpha\sin t+\cos t)/(1+\alpha^2)$ 核对。于是 $I(\alpha)=C-\arctan\alpha$。又由 $|\sin t|\leq t$，
\[
    |I(\alpha)|\leq\int_0^{+\infty}e^{-\alpha t}\,\mathrm{d}t
    =\frac1\alpha\longrightarrow0\qquad(\alpha\to+\infty),
\]
从而 $C=\pi/2$。利用零点处已经证明的右连续性，得
\begin{equation}
    \int_0^{+\infty}\frac{\sin t}{t}\,\mathrm{d}t
    =\lim_{\alpha\to0+}I(\alpha)=\frac\pi2.
    \label{eq:parameter-dirichlet-integral-value}
\end{equation}
这里的极限交换由一致收敛支持；仅算出 $\alpha>0$ 时的积分表达式，还不足以识别无阻尼积分。
\end{solve}

\subsection{含参变量瑕积分}
\label{subsec:parameter-dependent-singular-integrals}

设 $f\in C(E\times(c,d])$，其中 $c<d$ 为有限数，允许 $y=c$ 为瑕点。若每个 $x\in E$ 处都有有限极限
\[
    I(x)=\int_c^d f(x,y)\,\mathrm{d}y
    :=\lim_{\eta\to0+}\int_{c+\eta}^d f(x,y)\,\mathrm{d}y,
\]
则称其为\emph{含参变量瑕积分}（parameter-dependent improper integral with an endpoint singularity）。其一致收敛要求
\begin{equation}
    \forall\varepsilon>0\ \exists\eta_0\in(0,d-c)
    \forall\eta\in(0,\eta_0)\ \forall x\in E,
    \quad\left|\int_c^{c+\eta}f(x,y)\,\mathrm{d}y\right|<\varepsilon.
    \label{eq:parameter-singular-uniform-tail}
\end{equation}
等价的一致柯西条件为：对充分小的 $0<\eta_1<\eta_2$，$|\int_{c+\eta_1}^{c+\eta_2}f(x,y)\,\mathrm{d}y|$ 对全部 $x\in E$ 统一地小。

令 $t=1/(y-c)$，则 $y=c+1/t$，$\mathrm{d}y=-t^{-2}\mathrm{d}t$，从而
\begin{equation}
    \int_c^d f(x,y)\,\mathrm{d}y
    =\int_{1/(d-c)}^{+\infty}
        \frac{f(x,c+1/t)}{t^2}\,\mathrm{d}t.
    \label{eq:parameter-singular-to-infinite}
\end{equation}
两个截断误差完全对应，因此前面的一致柯西准则、魏尔斯特拉斯判别法、连续性、积分交换及积分号下求导定理均可转移过来。例如，若 $|f(x,y)|\leq g(y)$ 且非负函数 $g$ 在 $(c,d]$ 上瑕可积，就有绝对一致收敛；若 $f,f_x$ 联合连续、某个参数处原积分收敛，且 $\int_c^d f_x(x,y)\,\mathrm{d}y$ 内闭一致收敛，则 $I'(x)=\int_c^d f_x(x,y)\,\mathrm{d}y$。若两端都有瑕点，应固定一个内部点分别讨论两端，不能用同一个趋近参数使两端的发散部分互相抵消。

\begin{example}
\label{ex:parameter-power-singularity-uniformity}
讨论 $\displaystyle\int_0^1y^{s-1}\,\mathrm{d}y$ 的参数范围及一致收敛性。
\end{example}

\begin{solve}
当且仅当 $s>0$ 时积分收敛，值为 $1/s$。对任意 $\delta>0$、$0<\eta<1$ 及 $s\geq\delta$，
\[
    \int_0^\eta y^{s-1}\,\mathrm{d}y
    =\frac{\eta^s}{s}\leq\frac{\eta^\delta}{\delta}\longrightarrow0.
\]
因此在每个 $[\delta,+\infty)$ 上一致收敛。但在 $(0,+\infty)$ 上，固定任何 $\eta\in(0,1)$，$\eta^s/s\to+\infty$（$s\to0+$），故不一致收敛。参数逼近定义域边界时，可积奇性的强弱可能失去统一控制；这也解释了后面讨论 $\Gamma$ 与 $B$ 函数时为何把参数限制在正半轴的闭子区间上。
\end{solve}

含参变量积分还可推广到向量参数以及重积分、曲线积分和曲面积分。例如，设 $S$ 是固定的紧致、分片光滑、有指定侧的曲面，面积为 $|S|$，$\boldsymbol n$ 为其指定单位法向量；若 $\boldsymbol F(\boldsymbol r,t)$ 在 $S\times[a,b]$ 上连续，则瞬时有向通量
\[
    q(t)=\iint_S\boldsymbol F(\boldsymbol r,t)\cdot\boldsymbol n(\boldsymbol r)\,\mathrm{d}S
\]
连续，因为
\[
    |q(t)-q(t_0)|\leq |S|\sup_{\boldsymbol r\in S}
       |\boldsymbol F(\boldsymbol r,t)-\boldsymbol F(\boldsymbol r,t_0)|\longrightarrow0.
\]
这里即使分片交界处的法向量不连续，估计仍成立，因为法向量长度为 $1$，交界线不影响曲面积分。在流体模型中，若 $\boldsymbol F$ 表示体积通量密度，则 $[t_0,t_1]$ 内的净通过量为 $Q=\int_{t_0}^{t_1}q(t)\,\mathrm{d}t$；稳定场的 $q$ 与时间无关，因而 $Q=(t_1-t_0)q$。这一例子回到本章的动机，也说明关键在于积分区域上能够统一控制参数变化。

\section{\texorpdfstring{$\Gamma$}{Gamma} 函数与 \texorpdfstring{$B$}{B} 函数}
\label{sec:gamma-and-beta-functions}

前两节建立了由积分定义函数并研究其参数性质的方法。本节把它用于两类不能一般地用初等函数表示的积分：$\Gamma$ 函数把整数阶乘延伸到正实数参数，$B$ 函数统一处理两个端点的幂函数因子。先核实各自的收敛范围，再建立递推、光滑性与二者之间的关系，便能把一批定积分和几何量的计算归约到这些函数的特殊值。

\subsection{\texorpdfstring{$\Gamma$}{Gamma} 函数}
\label{subsec:gamma-function}

\begin{definition}[伽马函数]
\label{def:gamma-euler-integral}
对 $s>0$，定义\emph{伽马函数}（gamma function）
\begin{equation}
    \Gamma(s)=\int_0^{+\infty}t^{s-1}e^{-t}\,\mathrm{d}t.
    \label{eq:gamma-euler-integral}
\end{equation}
本章只讨论实参数的积分定义。
\end{definition}

在 $(0,1]$ 上，$e^{-1}t^{s-1}\leq t^{s-1}e^{-t}\leq t^{s-1}$，故零点附近当且仅当 $s>0$ 时可积；在 $[1,+\infty)$ 上，对任意固定实数 $s$，幂函数增长都可被指数衰减控制，因而无穷尾部收敛。因此式~\eqref{eq:gamma-euler-integral} 的实数定义域恰为 $(0,+\infty)$，且 $\Gamma(s)>0$。当 $0<s<1$ 时，它同时包含零点瑕积分与无穷积分，应在 $1$ 处分开理解。

\begin{proposition}[递推公式与阶乘]
\label{prop:gamma-recurrence-factorials}
对 $s>0$，有
\begin{equation}
    \Gamma(s+1)=s\Gamma(s),\qquad
    \Gamma(1)=1,\qquad \Gamma(n+1)=n!\quad(n=0,1,2,\ldots).
    \label{eq:gamma-recurrence-factorials}
\end{equation}
此外，换元 $t=u^2$ 得到
\begin{equation}
    \Gamma(s)=2\int_0^{+\infty}u^{2s-1}e^{-u^2}\,\mathrm{d}u,
    \qquad s>0.
    \label{eq:gamma-gaussian-form}
\end{equation}
\end{proposition}

\begin{proof}
先在 $[\eta,A]$ 上分部积分，再令 $\eta\to0+$、$A\to+\infty$。因为 $s>0$ 时 $t^se^{-t}$ 在两端均趋于零，故
\[
    \Gamma(s+1)=\bigl[-t^se^{-t}\bigr]_0^{+\infty}
       +s\int_0^{+\infty}t^{s-1}e^{-t}\,\mathrm{d}t
    =s\Gamma(s).
\]
由 $\Gamma(1)=\int_0^{+\infty}e^{-t}\,\mathrm{d}t=1$，归纳即得阶乘关系。平方换元也先在有限正常区间上进行，再取端点极限，得到式~\eqref{eq:gamma-gaussian-form}。
\end{proof}

\begin{theorem}[光滑性与各阶导数]
\label{thm:gamma-smoothness-derivatives}
$\Gamma\in C^\infty((0,+\infty))$，且对每个非负整数 $k$，
\begin{equation}
    \Gamma^{(k)}(s)=\int_0^{+\infty}
       t^{s-1}(\ln t)^k e^{-t}\,\mathrm{d}t,\qquad s>0.
    \label{eq:gamma-derivative-integrals}
\end{equation}
这些积分在每个闭区间 $[\delta,M]\subset(0,+\infty)$ 上绝对一致收敛。
\end{theorem}

\begin{proof}
固定 $0<\delta<M$ 及 $k\geq0$。对 $s\in[\delta,M]$，有
\[
    |t^{s-1}(\ln t)^k e^{-t}|
    \leq\begin{cases}
        t^{\delta-1}|\ln t|^k,&0<t\leq1,\\
        t^{M-1}(\ln t)^k e^{-t},&t\geq1.
    \end{cases}
\]
第一个控制函数经 $u=-\ln t$ 换元，其积分为 $\int_0^{+\infty}e^{-\delta u}u^k\,\mathrm{d}u<+\infty$；第二个控制函数在无穷远处可由常数倍的 $e^{-t/2}$ 控制，因而也可积。魏尔斯特拉斯判别法分别控制两端，保证每阶参数偏导数的积分绝对一致收敛。反复应用含参变量积分的求导定理及连续性定理，即得所有阶导数及其连续性。由于任意 $s>0$ 都包含在某个这样的闭区间内部，结论在整个正半轴成立。
\end{proof}

\begin{proposition}[严格凸性与严格对数凸性]
\label{prop:gamma-strict-log-convexity}
$\Gamma(s)$ 与 $\ln\Gamma(s)$ 都是 $(0,+\infty)$ 上的严格凸函数。
\end{proposition}

\begin{proof}
由式~\eqref{eq:gamma-derivative-integrals}，
\[
    \Gamma''(s)=\int_0^{+\infty}t^{s-1}(\ln t)^2e^{-t}\,\mathrm{d}t>0,
\]
因为被积函数非负，且在不包含 $1$ 的任一非退化正区间上为正。又因 $\Gamma>0$，
\[
    (\ln\Gamma(s))''
    =\frac{\Gamma(s)\Gamma''(s)-(\Gamma'(s))^2}{\Gamma(s)^2}.
\]
对 $u(t)=t^{(s-1)/2}e^{-t/2}$ 与 $v(t)=u(t)\ln t$ 应用积分形式的柯西--施瓦茨不等式，得 $(\Gamma'(s))^2\leq\Gamma(s)\Gamma''(s)$。为核查严格性，令 $m=\Gamma'(s)/\Gamma(s)$，则
\[
    \Gamma''(s)-\frac{(\Gamma'(s))^2}{\Gamma(s)}
    =\int_0^{+\infty}t^{s-1}e^{-t}(\ln t-m)^2\,\mathrm{d}t>0.
\]
右端的连续非负被积函数只在 $t=e^m$ 处为零，在其他点附近均给出正积分，故 $(\ln\Gamma)''>0$。
\end{proof}

\subsection{\texorpdfstring{$B$}{B} 函数}
\label{subsec:beta-function}

\begin{definition}[贝塔函数]
\label{def:beta-euler-integral}
对 $p,q>0$，定义\emph{贝塔函数}（beta function）
\begin{equation}
    B(p,q)=\int_0^1 t^{p-1}(1-t)^{q-1}\,\mathrm{d}t.
    \label{eq:beta-euler-integral}
\end{equation}
它是含两个参数的广义积分。
\end{definition}

在零点附近，$(1-t)^{q-1}$ 是有正下界的连续函数，故可积性要求 $p>0$；在 $1$ 附近同理要求 $q>0$。所以该积分的实数定义域恰为 $(0,+\infty)^2$，且在此域内恒正。固定任意内部分点，例如 $1/2$，分别取两个端点的瑕积分极限，便得到严格的定义；当某端点没有奇性时，该部分就是普通定积分。

\begin{proposition}[对称性与递推公式]
\label{prop:beta-symmetry-recurrence}
对 $p,q>0$，有
\begin{align}
    B(p,q)&=B(q,p),\label{eq:beta-symmetry}\\
    B(p+1,q)&=\frac{p}{p+q}B(p,q),\qquad
    B(p,q+1)=\frac{q}{p+q}B(p,q).
    \label{eq:beta-forward-recurrence}
\end{align}
特别地，
\begin{align}
    B(p,q)&=\frac{p-1}{p+q-1}B(p-1,q)
        &&(p>1,\ q>0),\label{eq:beta-backward-recurrence}\\
    B(p,q)&=\frac{(p-1)(q-1)}{(p+q-1)(p+q-2)}B(p-1,q-1)
        &&(p>1,\ q>1).\label{eq:beta-double-recurrence}
\end{align}
\end{proposition}

\begin{proof}
在定义式中令 $t=1-u$，即得对称性。由 $1=t+(1-t)$，有
\[
    B(p,q)=B(p+1,q)+B(p,q+1).
\]
另一方面，$t^p(1-t)^q$ 在两端均趋于零，其导数的两个分项在 $p,q>0$ 时分别可积。先在 $[\eta,1-\eta]$ 上积分导数再令 $\eta\to0+$，得到
\[
    0=pB(p,q+1)-qB(p+1,q).
\]
联立即得向前递推式；对参数平移并重复使用，即得后两式。向后递推时必须保留 $p>1$、$q>1$ 等条件，保证右侧的积分仍在定义域内。
\end{proof}

\begin{proposition}[三角形式与无穷积分形式]
\label{prop:beta-alternative-integrals}
对 $p,q>0$，有
\begin{align}
    B(p,q)&=2\int_0^{\pi/2}
       (\cos\theta)^{2p-1}(\sin\theta)^{2q-1}\,\mathrm{d}\theta,
       \label{eq:beta-trigonometric-form}\\
    B(p,q)&=\int_0^{+\infty}
       \frac{u^{p-1}}{(1+u)^{p+q}}\,\mathrm{d}u.
       \label{eq:beta-infinite-form}
\end{align}
\end{proposition}

\begin{proof}
第一式在定义式中取 $t=\cos^2\theta$，注意积分方向反转；第二式取 $t=u/(1+u)$，此时 $\mathrm{d}t=(1+u)^{-2}\mathrm{d}u$。两次换元先在避开端点的闭区间上进行，再分别取端点极限。由对称性，式~\eqref{eq:beta-trigonometric-form} 中的正弦、余弦幂次也可以互换。
\end{proof}

\subsection{\texorpdfstring{$\Gamma$}{Gamma} 函数与 \texorpdfstring{$B$}{B} 函数的关系}
\label{subsec:relation-between-gamma-and-beta-functions}

\begin{theorem}[贝塔--伽马关系]
\label{thm:beta-gamma-relation}
对 $p,q>0$，有
\begin{equation}
    B(p,q)=\frac{\Gamma(p)\Gamma(q)}{\Gamma(p+q)}.
    \label{eq:beta-gamma-relation}
\end{equation}
从而 $B$ 在 $(0,+\infty)^2$ 内具有任意阶连续偏导数。
\end{theorem}

\begin{proof}
由式~\eqref{eq:gamma-gaussian-form}，两个一维正积分的乘积可以写成第一象限上的广义二重积分
\[
    \Gamma(p)\Gamma(q)
    =4\iint_{x>0,\,y>0}x^{2p-1}y^{2q-1}e^{-(x^2+y^2)}
       \,\mathrm{d}x\,\mathrm{d}y.
\]
为说明换元的合法性，先在矩形 $[\varepsilon,R]^2$ 上积分，再令 $\varepsilon\downarrow0$、$R\uparrow+\infty$，矩形积分等于两个截断一维积分的乘积，故极限存在。被积函数非负，也可改用极坐标截断区域 $\delta\leq r\leq L$、$\eta\leq\theta\leq\pi/2-\eta$：每个这样的区域包含于某个矩形截断区域，每个矩形截断区域也包含于某个这样的区域，所以两族积分的上确界相同。这为改换穷竭区域以及随后的极坐标变换提供依据，即使坐标轴附近有瑕点也成立。

令 $x=r\cos\theta$、$y=r\sin\theta$，雅可比因子为 $r$，于是
\begin{align*}
    \Gamma(p)\Gamma(q)
    &=4\left(\int_0^{+\infty}r^{2(p+q)-1}e^{-r^2}\,\mathrm{d}r\right)
       \left(\int_0^{\pi/2}\cos^{2p-1}\theta\sin^{2q-1}\theta
          \,\mathrm{d}\theta\right)\\
    &=\Gamma(p+q)B(p,q).
\end{align*}
两个因子都收敛，且 $\Gamma(p+q)>0$，故可相除。最后，$\Gamma$ 光滑且在正半轴不为零，关系式便给出 $B$ 的光滑性。
\end{proof}

\begin{theorem}[余元公式]
\label{thm:gamma-reflection-formula}
若 $0<p<1$，则
\begin{equation}
    \Gamma(p)\Gamma(1-p)=B(p,1-p)=\frac{\pi}{\sin(\pi p)}.
    \label{eq:gamma-reflection-formula}
\end{equation}
该恒等式也称\emph{反射公式}（reflection formula）；这里的参数范围保证两边的 $\Gamma$ 函数都由本章的正参数积分定义。
\end{theorem}

\begin{proof}
贝塔--伽马关系与 $\Gamma(1)=1$ 给出第一个等号。由式~\eqref{eq:beta-infinite-form}，在 $[1,+\infty)$ 部分取 $u=1/t$，得
\[
    B(p,1-p)=\int_0^{+\infty}\frac{t^{p-1}}{1+t}\,\mathrm{d}t
    =\int_0^1\frac{t^{p-1}+t^{-p}}{1+t}\,\mathrm{d}t.
\]
使用有限几何恒等式
\[
    \frac1{1+t}=\sum_{k=0}^N(-1)^kt^k
       +\frac{(-1)^{N+1}t^{N+1}}{1+t}.
\]
余项乘以 $t^{p-1}+t^{-p}$ 后的积分绝对值不超过
$1/(N+p+1)+1/(N+2-p)$，故趋于零，无须假定几何级数在 $[0,1]$ 上一致收敛。逐项计算有限和并在极限中合并相邻指标，得到
\begin{equation}
    B(p,1-p)
    =\lim_{N\to\infty}\sum_{k=0}^N(-1)^k
       \left(\frac1{k+p}+\frac1{k+1-p}\right)
    =\frac1p+\sum_{k=1}^{\infty}(-1)^k\frac{2p}{p^2-k^2}.
    \label{eq:beta-reflection-series}
\end{equation}
合并时只差一个趋于零的末项；合并后的级数绝对收敛。

下面按教材思路用 $\cos(px)$ 的傅里叶系数计算这个级数，并补出在 $x=0$ 处取值的依据。直接积分得
\[
    \frac1\pi\int_0^\pi\cos(pt)\,\mathrm{d}t
       =\frac{\sin(\pi p)}{\pi p},\qquad
    \frac2\pi\int_0^\pi\cos(pt)\cos(kt)\,\mathrm{d}t
       =\frac{\sin(\pi p)}\pi\frac{2p(-1)^k}{p^2-k^2}.
\]
因此对应傅里叶部分和在零点的值为
\[
    S_N=\frac{\sin(\pi p)}\pi
       \left(\frac1p+\sum_{k=1}^N\frac{2p(-1)^k}{p^2-k^2}\right).
\]
利用有限三角恒等式 $1+2\sum_{k=1}^N\cos(kt)=\sin((N+1/2)t)/\sin(t/2)$，以及常数函数的相应积分为 $1$，有
\[
    S_N-1=\frac1\pi\int_0^\pi
      \frac{\cos(pt)-1}{\sin(t/2)}\sin((N+1/2)t)\,\mathrm{d}t.
\]
函数 $h(t)=(\cos(pt)-1)/\sin(t/2)$ 在 $t=0$ 处补为零后属于 $C^1([0,\pi])$。对右侧分部积分，其绝对值不超过
\[
    \frac{|h(0)|+|h(\pi)|+\int_0^\pi|h'(t)|\,\mathrm{d}t}
         {\pi(N+1/2)}\longrightarrow0.
\]
故 $S_N\to1$。结合式~\eqref{eq:beta-reflection-series} 即得 $B(p,1-p)=\pi/\sin(\pi p)$。
\end{proof}

取 $p=1/2$，由 $\Gamma(1/2)>0$ 得
\begin{equation}
    \Gamma\left(\frac12\right)=\sqrt\pi,\qquad
    \Gamma\left(n+\frac12\right)
       =\frac{(2n)!}{4^n n!}\sqrt\pi\quad(n=0,1,2,\ldots).
    \label{eq:gamma-half-integer-values}
\end{equation}
其中第二式由递推公式得到。再由式~\eqref{eq:gamma-gaussian-form}，得到高斯积分
\begin{equation}
    \int_0^{+\infty}e^{-t^2}\,\mathrm{d}t=\frac{\sqrt\pi}{2},
    \qquad\int_{-\infty}^{+\infty}e^{-t^2}\,\mathrm{d}t=\sqrt\pi.
    \label{eq:gamma-gaussian-integral}
\end{equation}
递推公式将正参数处的值归约到 $(0,1]$ 上的值，余元公式又联系 $p$ 与 $1-p$；这些公式的使用都应遵守各自的参数范围。

\begin{example}
\label{ex:beta-sine-power-integrals}
计算 $\displaystyle I_n=\int_0^{\pi/2}\sin^n x\,\mathrm{d}x$（$n=0,1,2,\ldots$）。
\end{example}

\begin{solve}
由贝塔函数的三角形式及其与伽马函数的关系，
\begin{equation}
    I_n=\frac12B\left(\frac12,\frac{n+1}{2}\right)
       =\frac{\sqrt\pi\,\Gamma((n+1)/2)}{2\Gamma((n+2)/2)}
       =\begin{cases}
           \dfrac{(n-1)!!}{n!!}\dfrac\pi2,&n\text{ 为偶数},\\
           \dfrac{(n-1)!!}{n!!},&n\text{ 为奇数}.
         \end{cases}
    \label{eq:beta-sine-power-integrals}
\end{equation}
这里 $n!!=n(n-2)\cdots$，约定 $0!!=(-1)!!=1$；故 $n=0$ 时仍给出 $I_0=\pi/2$。换元 $x\mapsto\pi/2-x$ 可知余弦同次幂的积分相等。中间的伽马函数表达式还对任意实数 $n>-1$ 成立，此范围正是端点处可积的条件。
\end{solve}

\begin{example}
\label{ex:gamma-unit-ball-volume}
求 $\mathbb R^n$ 中单位球体 $D_n=\{(x_1,\ldots,x_n):x_1^2+\cdots+x_n^2\leq1\}$ 的体积 $V_n$。
\end{example}

\begin{solve}
先用截面递推使计算中的维数关系清楚。$V_1=2$；$n\geq2$ 时，在 $x_n=t$ 处的截面是半径为 $\sqrt{1-t^2}$ 的 $(n-1)$ 维球体，其体积由缩放变换等于 $V_{n-1}(1-t^2)^{(n-1)/2}$。球面作为有限个局部光滑图形之并具有零 $n$ 维体积，因而球体的示性函数可积，截面积分公式适用。由 $u=t^2$ 换元，
\begin{align*}
    V_n&=V_{n-1}\int_{-1}^1(1-t^2)^{(n-1)/2}\,\mathrm{d}t\\
       &=V_{n-1}B\left(\frac12,\frac{n+1}{2}\right)
        =V_{n-1}\frac{\sqrt\pi\,\Gamma((n+1)/2)}{\Gamma((n+2)/2)}.
\end{align*}
由归纳法及 $\Gamma(3/2)=\sqrt\pi/2$，得到
\begin{equation}
    V_n=\frac{\pi^{n/2}}{\Gamma(n/2+1)}
    =\begin{cases}
       \dfrac{\pi^m}{m!},&n=2m,\ m\geq1,\\
       \dfrac{2^{m+1}\pi^m}{(2m+1)!!},&n=2m+1,\ m\geq0.
     \end{cases}
    \label{eq:gamma-unit-ball-volume}
\end{equation}
例如 $V_2=\pi$、$V_3=4\pi/3$；半径为 $R>0$ 的球体体积再乘 $R^n$。

教材使用的高维球坐标给出同一结果：在 $n\geq2$ 时，径向变量 $r\in(0,1)$，前 $n-2$ 个角属于 $(0,\pi)$，最后一个角属于 $(0,2\pi)$；逐次作平面极坐标变换，雅可比绝对值为
\[
    r^{n-1}\sin^{n-2}\theta_1\sin^{n-3}\theta_2\cdots\sin\theta_{n-2}.
\]
坐标退化集及角度分界集为零体积集，不影响积分。于是
\[
    V_n=\frac{2\pi}{n}\prod_{j=1}^{n-2}
         \int_0^\pi\sin^j\theta\,\mathrm{d}\theta
       =\frac{2\pi}{n}\prod_{j=1}^{n-2}
         \frac{\sqrt\pi\,\Gamma((j+1)/2)}{\Gamma((j+2)/2)}
       =\frac{\pi^{n/2}}{\Gamma(n/2+1)}.
\]
$n=2$ 时乘积为空，按 $1$ 计算；一般情形的伽马函数因子逐项消去，与截面递推结果一致。
\end{solve}

\section*{习题十七}
\addcontentsline{toc}{section}{习题十七}
